%=============================================================================
% FORMAL PROOF: Branch-by-Branch Entanglement of \psi from
% Monotone-Operator Uniqueness
%
% Shows that the branch-by-branch picture for quantum matter coupled
% to the DFD scalar field follows from the uniqueness theorems already
% proven in Section III / Appendix U of [DFDUnified].
%=============================================================================

\documentclass[aps,prd,twocolumn,superscriptaddress,nofootinbib]{revtex4-2}
\usepackage{amsmath,amssymb,amsthm}
\usepackage{braket}
\usepackage{tcolorbox}

\newtheorem{theorem}{Theorem}
\newtheorem{lemma}[theorem]{Lemma}
\newtheorem{proposition}[theorem]{Proposition}
\newtheorem{corollary}[theorem]{Corollary}
\newtheorem{definition}[theorem]{Definition}
\newtheorem*{remark}{Remark}

\newcommand{\Hm}{\mathcal{H}_{\mathrm{matter}}}
\newcommand{\Hpsi}{\mathcal{H}_{\psi}}
\newcommand{\astar}{a_\star}
\newcommand{\mufunction}{\mu}

\begin{document}

\title{Branch-by-Branch Entanglement of the DFD Scalar Field:\\
A Theorem from Monotone-Operator Uniqueness}

\author{Research Note (DFD Cross-Reference Series)}
\date{\today}

\maketitle

%=============================================================================
\section{Summary of the Result}
%=============================================================================

We prove that the branch-by-branch entanglement structure
\begin{equation}\label{eq:main-result}
\ket{\Psi_{\mathrm{total}}} = \sum_i c_i \ket{i}_{\mathrm{matter}} \otimes
\ket{\psi[\rho_i]}_\psi
\end{equation}
follows from exactly two inputs:
\begin{enumerate}
\item The \textbf{linearity of quantum mechanics} (standard, unmodified).
\item The \textbf{uniqueness of the static DFD solution} $\psi$ given a
source $\rho$, as established by Theorem~U.4 (= Theorem~3.2) of
\cite{DFDUnified} via the Browder--Minty theorem for strictly monotone
operators.
\end{enumerate}

No additional postulate is required. The result is a \emph{logical
consequence} of the existing mathematical framework.

%=============================================================================
\section{The Uniqueness Theorem (Extracted from [DFDUnified])}
%=============================================================================

We state the exact theorem from Appendix~U of \cite{DFDUnified}.

\begin{tcolorbox}[colback=blue!5!white, colframe=blue!60!black,
title=Theorem U.4 (Static Uniqueness)]
\textbf{Setting.} Let $\Omega \subseteq \mathbb{R}^3$ be a domain
(bounded or exterior with asymptotic flatness). Define the flux operator
$\mathbf{a}(\xi) := \mufunction(|\xi|)\xi$, where
$\mufunction: [0,\infty) \to (0,\infty)$ satisfies:
\begin{itemize}
\item[\textbf{(A1)}] Continuity on $[0,\infty)$.
\item[\textbf{(A2)}] Coercivity: $\mufunction(|\xi|)|\xi|^2 \geq \alpha|\xi|^p$
for some $\alpha > 0$, $p \geq 2$.
\item[\textbf{(A3)}] Growth bound: $|\mufunction(|\xi|)\xi| \leq \beta(1+|\xi|)^{p-1}$.
\item[\textbf{(A4)}] Strict monotonicity:
$(\mathbf{a}(\xi) - \mathbf{a}(\eta))\cdot(\xi - \eta) > 0$ for $\xi \neq \eta$.
\end{itemize}

\textbf{Statement.} For any source $f \in (W_0^{1,p}(\Omega))'$, the weak
problem
\begin{equation}
\int_\Omega \mathbf{a}(\nabla\psi)\cdot\nabla v\,d^3x =
\int_\Omega f\,v\,d^3x, \quad \forall\,v \in W_0^{1,p}(\Omega),
\end{equation}
with prescribed Dirichlet boundary data, has \textbf{exactly one} solution
$\psi \in W^{1,p}(\Omega)$.

\textbf{Proof method.} Existence: Browder--Minty theorem (monotone, coercive,
hemicontinuous operator). Uniqueness: strict convexity of the energy functional
$\mathcal{E}[\psi] = \int H(\nabla\psi)\,dx - \int f\psi\,dx$.
\end{tcolorbox}

\begin{remark}[On ``Leray--Schauder'' vs.\ Browder--Minty]
The user's original argument referenced the Leray--Schauder fixed-point
theorem. The DFD paper actually uses the \textbf{Browder--Minty theorem}
(for monotone operators in reflexive Banach spaces), which is a more direct
and sharper tool for this class of quasilinear elliptic PDEs. Both yield
uniqueness, but Browder--Minty does so under weaker hypotheses. The logical
structure of the branch-by-branch proof is identical for either method;
what matters is that $\rho \mapsto \psi[\rho]$ is a well-defined function.
\end{remark}

\begin{lemma}[DFD $\mufunction$ satisfies strict monotonicity]\label{lem:strict-mono}
For $\mufunction(x) = x/(1+x)$, the flux $\mathbf{a}(\xi) = \mufunction(|\xi|)\xi$
is strictly monotone.
\end{lemma}
\begin{proof}
Compute $\partial a_i / \partial \xi_j = \mufunction(|\xi|)\delta_{ij} +
\mufunction'(|\xi|)\xi_i\xi_j/|\xi|$. Since $\mufunction'(s) = 1/(1+s)^2 > 0$
for all $s \geq 0$, the Jacobian matrix is positive definite (not merely
semi-definite). Hence strict monotonicity holds. (This is Lemma~U.1 of
\cite{DFDUnified}.)
\end{proof}

%=============================================================================
\section{Verification of Hypotheses for Each Branch}
%=============================================================================

We now verify that the uniqueness theorem applies separately to each
quantum branch.

\begin{proposition}[Branch applicability]\label{prop:branch-apply}
Let $\ket{\Psi} = \sum_i c_i \ket{i}$ be a quantum state of matter,
where each $\ket{i}$ is an eigenstate of the mass-energy density operator
with eigenvalue $\rho_i(\mathbf{x})$. Suppose each $\rho_i$ satisfies:
\begin{enumerate}
\item[(i)] $\rho_i \geq 0$ on $\Omega$,
\item[(ii)] $\rho_i \in L^q(\Omega)$ with $q > 3/p'$ (so that $f_i = (8\pi G/c^2)\rho_i$
lies in $(W_0^{1,p}(\Omega))'$),
\item[(iii)] $\rho_i$ is localized (compact support or sufficient decay at infinity).
\end{enumerate}
Then the hypotheses of Theorem~U.4 are satisfied for each source $f_i$,
independently of the other branches.
\end{proposition}

\begin{proof}
The hypotheses (A1)--(A4) concern the flux operator $\mathbf{a}$ and the
function $\mufunction$, not the source. They are verified once and for all
in Lemma~\ref{lem:strict-mono}. The only requirement on the source is that
$f_i \in (W_0^{1,p}(\Omega))'$, which is guaranteed by condition~(ii).
Conditions (i) and (iii) are physically natural for each branch but are
\emph{not required} by the theorem---the uniqueness holds for any
$f \in V'$, including signed or oscillatory sources.

Crucially, the theorem places \textbf{no requirement} that $\rho$ be
``classical'' or ``definite'' in any ontological sense. It requires only
that $\rho$ be an element of a function space. Each branch density
$\rho_i$ is such an element.
\end{proof}

\begin{remark}[What could break?]
The hypotheses that could potentially fail in exotic quantum scenarios:
\begin{itemize}
\item \textbf{Point particles:} If $\rho_i = m\,\delta^3(\mathbf{x} - \mathbf{x}_0)$,
then $f_i \notin L^q$ for any $q > 1$. However, $\delta^3 \in (W_0^{1,p}(\Omega))'$
for $p > 3$ by Sobolev embedding, so existence holds for $p > 3$. The DFD $\mufunction$
has $p = 2$, so point-particle sources require regularization (e.g., smearing over a
Compton wavelength). This is a standard issue shared with Newtonian gravity and GR.
\item \textbf{Unbounded domain:} For $\Omega = \mathbb{R}^3$, one uses weighted
Sobolev spaces. Theorem~U.7 (Exterior Well-Posedness) handles this case.
\item \textbf{Non-static sources:} The theorem is for the static equation.
Time-dependent sources require the hyperbolic theory (Theorem~U.8).
\end{itemize}
\end{remark}

%=============================================================================
\section{The Main Theorem: Branch-by-Branch Entanglement}
%=============================================================================

\begin{definition}[Solution map]
The \textbf{DFD solution map} is the function
\begin{equation}
\mathcal{S}: (W_0^{1,p}(\Omega))' \to W^{1,p}(\Omega), \quad
f \mapsto \psi[f],
\end{equation}
where $\psi[f]$ is the unique weak solution of the static DFD equation
with source $f$ and fixed boundary data. This map is well-defined by
Theorems~U.3 (existence) and U.4 (uniqueness).
\end{definition}

\begin{theorem}[Branch-by-Branch Entanglement]\label{thm:branch}
Let the following hold:
\begin{enumerate}
\item[\textbf{(QM)}] Quantum mechanics is linear: $\ket{\Psi} = \sum_i c_i \ket{i}$
with $\sum_i |c_i|^2 = 1$.
\item[\textbf{(Def)}] Each $\ket{i}$ has a definite mass-energy density
$\rho_i \in L^q(\Omega)$ with $q > 3/p'$.
\item[\textbf{(QS)}] The quasi-static approximation holds: the field $\psi$
adjusts instantaneously to the source distribution (valid when
matter evolution timescale $\gg \Omega/c$).
\item[\textbf{(Uniq)}] The DFD static field equation has a unique solution
for each source (Theorem~U.4 of \cite{DFDUnified}).
\end{enumerate}

Then the total quantum state in the extended Hilbert space
$\Hm \otimes \Hpsi$ is:
\begin{equation}\label{eq:branch-state}
\boxed{
\ket{\Psi_{\mathrm{total}}} = \sum_i c_i \ket{i}_{\mathrm{matter}}
\otimes \ket{\psi[\rho_i]}_\psi
}
\end{equation}
where $\psi[\rho_i] := \mathcal{S}[(8\pi G/c^2)\rho_i]$ is the unique
DFD solution sourced by $\rho_i$.
\end{theorem}

\begin{proof}
The proof proceeds in three steps.

\medskip
\noindent\textbf{Step 1: The solution map is a function.}
By \textbf{(Uniq)} (Theorem~U.4), the map $\rho \mapsto \psi[\rho]$
assigns to each source density a \emph{unique} field configuration.
That is, $\mathcal{S}$ is a well-defined function, not a
multi-valued correspondence.

\medskip
\noindent\textbf{Step 2: Functional constraint + superposition $\Rightarrow$ entanglement.}
By \textbf{(Def)}, each basis state $\ket{i}$ has a definite density $\rho_i$.
By Step~1, each $\rho_i$ determines a unique $\psi_i := \psi[\rho_i]$.
The physical constraint (the field equation) therefore correlates each
matter branch with a definite field configuration:
\begin{equation}
\ket{i}_{\mathrm{matter}} \longleftrightarrow \ket{\psi_i}_\psi.
\end{equation}

By \textbf{(QM)} (linearity), the total state is a superposition of
these correlated states:
\begin{equation}
\ket{\Psi_{\mathrm{total}}} = \sum_i c_i \ket{i, \psi_i}.
\end{equation}

This is the standard quantum-mechanical rule: if a system in state
$\sum c_i \ket{i}$ is subjected to a \emph{deterministic} constraint
that maps $\ket{i} \to \ket{i} \otimes \ket{\phi_i}$, the result is
$\sum c_i \ket{i}\otimes\ket{\phi_i}$. The determinism comes from
uniqueness (Step~1), and the linearity is \textbf{(QM)}.

\medskip
\noindent\textbf{Step 3: The quasi-static limit.}
By \textbf{(QS)}, the field equation is solved as a static boundary
value problem for each instantaneous matter configuration. This is the
standard Born--Oppenheimer-type approximation: the ``fast'' degree of
freedom ($\psi$) adjusts adiabatically to the ``slow'' one (matter).
The result~\eqref{eq:branch-state} is the adiabatic entangled state.
\end{proof}

%=============================================================================
\section{Contrast with Expectation-Value Sourcing}
%=============================================================================

\begin{proposition}[Expectation-value sourcing gives a different $\psi$]\label{prop:nonlinear}
Define the expectation-value source:
\begin{equation}
\rho_{\mathrm{avg}} := \sum_i |c_i|^2 \rho_i = \braket{\Psi|\hat{\rho}|\Psi}.
\end{equation}
Then generically:
\begin{equation}
\psi[\rho_{\mathrm{avg}}] \neq \sum_i |c_i|^2 \psi[\rho_i].
\end{equation}
\end{proposition}

\begin{proof}
The solution map $\mathcal{S}$ is \textbf{nonlinear} because the field
equation $-\nabla\cdot(\mufunction(|\nabla\psi|)\nabla\psi) = f$ is
nonlinear in $\psi$. Specifically, $\mufunction(|\nabla\psi|)$ depends
on $|\nabla\psi|$, so the superposition principle fails for the PDE:
\begin{equation}
\mathcal{S}[\lambda f_1 + (1-\lambda)f_2] \neq
\lambda\,\mathcal{S}[f_1] + (1-\lambda)\mathcal{S}[f_2]
\end{equation}
for generic $f_1 \neq f_2$ and $\lambda \in (0,1)$.

Therefore the two prescriptions---branch-by-branch sourcing
(Eq.~\ref{eq:branch-state}) and expectation-value sourcing---yield
\textbf{different physical predictions}. The nonlinearity of $\mufunction$
is what makes the distinction observable in principle.
\end{proof}

\begin{remark}[Physical consequence]
This is analogous to the Schr\"odinger--Newton debate in semiclassical
gravity: should one source the gravitational field with $\braket{\hat{T}_{\mu\nu}}$
(semiclassical, = expectation-value sourcing) or branch-by-branch? In DFD,
the uniqueness theorem \emph{forces} the branch-by-branch picture as the
logically consistent one, given the linearity of QM. The expectation-value
prescription would require abandoning either QM linearity or the uniqueness
of $\psi[\rho]$.
\end{remark}

%=============================================================================
\section{Precise Statement of What Is Proven}
%=============================================================================

\subsection{What this proof establishes}

\begin{enumerate}
\item \textbf{In the quasi-static limit} (no time derivatives in the
field equation), the branch-by-branch entangled state
$\ket{\Psi_{\mathrm{total}}} = \sum c_i \ket{i, \psi[\rho_i]}$ follows
from QM linearity + DFD uniqueness. \textbf{No additional postulate
is needed.}

\item Each branch density $\rho_i$ must be an element of
$(W_0^{1,p}(\Omega))'$ (i.e., a distribution no worse than a Sobolev
dual element). This is satisfied by any $L^q$ density with
$q > 3/p' = 3$ (for $p = 2$), and more broadly by any finite-energy
matter distribution.

\item The proof uses the Browder--Minty theorem (monotone operator theory),
which is the method actually employed in \cite{DFDUnified}. The same
conclusion follows from Leray--Schauder if one prefers that route, since
both yield the key ingredient: \textbf{uniqueness of $\psi$ given $\rho$}.

\item The continuous dependence (Theorem~U.9 of \cite{DFDUnified},
stability estimate) ensures that \emph{nearby} branch densities produce
\emph{nearby} field configurations. This gives the entanglement structure
a form of continuity: small perturbations of $\rho_i$ produce small
perturbations of $\psi[\rho_i]$.
\end{enumerate}

\subsection{What this proof does NOT establish}

\begin{enumerate}
\item \textbf{The dynamical (time-dependent) case.} The hyperbolic
evolution (Theorem~U.8) also has local uniqueness, so a time-dependent
version of the branch-by-branch picture should hold locally in time.
However, the quasi-static approximation is essential for the
\emph{instantaneous} entanglement structure~\eqref{eq:branch-state}.
The fully dynamical case requires a careful treatment of retardation
effects and is left open.

\item \textbf{Point-particle sources.} For $\rho_i = m\,\delta^3(\mathbf{x})$,
the source is too singular for the $p = 2$ Sobolev framework. Physical
regularization (smearing over a Compton wavelength or using the quantum
wave packet as the effective source) resolves this. The proof applies to
the regularized source.

\item \textbf{Decoherence corrections.} The entangled state~\eqref{eq:branch-state}
is an exact superposition. In practice, the different $\psi[\rho_i]$
carry distinguishable gravitational information, which may lead to
environment-induced decoherence. The rate and mechanism of such
decoherence is a separate question not addressed here.

\item \textbf{Full quantum gravity.} This proof treats $\psi$ as a
classical field entangled with quantum matter (the semi-classical regime).
Full quantization of $\psi$ is not needed: as noted in \cite{DFDUnified}
\S3.6, the action scales as $S_\psi \sim (M_P/\astar)^2 \gg \hbar$,
so quantum fluctuations of $\psi$ itself are negligible.

\item \textbf{Whether decoherence from $\psi$-branching is fast enough
to explain the quantum-to-classical transition.} This is an important
physical question but requires estimating the overlap
$\braket{\psi[\rho_i]|\psi[\rho_j]}$ for macroscopically distinct
$\rho_i, \rho_j$, which is a separate calculation.
\end{enumerate}

%=============================================================================
\section{Logical Structure (Schematic)}
%=============================================================================

The chain of logic is:

\begin{equation*}
\boxed{
\begin{aligned}
&\text{Browder--Minty (Thm U.3--U.4)} \\
&\quad \Downarrow \\
&\rho \mapsto \psi[\rho] \text{ is a well-defined \emph{function}} \\
&\quad \Downarrow \\
&\text{Each branch } \ket{i} \text{ with } \rho_i \text{ determines unique } \psi_i \\
&\quad \Downarrow \quad (\text{QM linearity}) \\
&\ket{\Psi_{\mathrm{total}}} = \sum c_i \ket{i} \otimes \ket{\psi_i} \\
&\quad \Downarrow \quad (\text{nonlinearity of } \mufunction) \\
&\psi[\langle\rho\rangle] \neq \langle\psi[\rho]\rangle
\;\Rightarrow\; \text{distinct from semiclassical sourcing}
\end{aligned}
}
\end{equation*}

%=============================================================================
\section{Technical Note: Browder--Minty vs.\ Leray--Schauder}
%=============================================================================

The DFD paper uses the \textbf{Browder--Minty theorem} (1963), which states
that a monotone, coercive, hemicontinuous operator from a reflexive Banach
space to its dual is surjective. Combined with strict monotonicity, this
gives uniqueness. This is the natural tool for the variational structure of
quasilinear elliptic PDEs.

The \textbf{Leray--Schauder theorem} (1934) is a topological fixed-point
theorem based on degree theory. It can also be used to prove existence for
the DFD equation, but requires additional a~priori estimates (to establish
compactness of the solution operator). The two approaches are complementary:

\begin{center}
\begin{tabular}{lcc}
\hline
Property & Browder--Minty & Leray--Schauder \\
\hline
Existence & \checkmark & \checkmark \\
Uniqueness & \checkmark (strict mono.) & Requires separate argument \\
Method & Variational/operator & Topological degree \\
Hypotheses & (A1)--(A4) & A priori bounds + compactness \\
Used in [DFDUnified] & \checkmark & Not explicitly \\
\hline
\end{tabular}
\end{center}

For the branch-by-branch proof, the critical ingredient is \textbf{uniqueness},
which Browder--Minty provides directly. If one uses Leray--Schauder for
existence, uniqueness still follows from the strict convexity of the energy
functional (an independent argument also given in [DFDUnified]).

Either way, the conclusion is the same: $\rho \mapsto \psi[\rho]$ is a
function, and the branch-by-branch picture follows.

\begin{thebibliography}{9}
\bibitem{DFDUnified} G.~T.~Alcock, ``Density Field Dynamics: A Complete
Unified Theory,'' v108 (2025). Section~III and Appendix~U.
\bibitem{Brezis} H.~Br\'ezis, \textit{Functional Analysis, Sobolev Spaces
and Partial Differential Equations}, Springer (2011).
\bibitem{Evans} L.~C.~Evans, \textit{Partial Differential Equations},
AMS Graduate Studies in Mathematics (2010).
\bibitem{GilbargTrudinger} D.~Gilbarg and N.~S.~Trudinger,
\textit{Elliptic Partial Differential Equations of Second Order},
Springer (2001).
\bibitem{LadyzhenskayaUraltseva} O.~A.~Ladyzhenskaya and N.~N.~Ural'tseva,
\textit{Linear and Quasilinear Elliptic Equations}, Academic Press (1968).
\bibitem{CrandallLiggett} M.~G.~Crandall and T.~M.~Liggett,
``Generation of semi-groups of nonlinear transformations on general
Banach spaces,'' \textit{Amer.~J.~Math.} \textbf{93}, 265--298 (1971).
\end{thebibliography}

\end{document}
